% STATUS: written 2026-11-06 (English gate clean; sources read in full as recorded below).
% Sources read in full: Leibniz 1684 (Acta Eruditorum, October 1684, 467--473; scan archive.org s1id13206500, rules of the calculus on p. 467 checked in the OCR text), Jacobson 1937 (Amer. J. Math. 59, 205--220; the definition (2c) on p. 206).
\section[The Leibniz rule: Leibniz 1684][The Leibniz rule]{The Leibniz rule: Leibniz 1684}\label{sec:origins:leibniz-rule}

Every derivation in this book satisfies the same identity:
\[
  D(xy) = D(x)\,y + x\,D(y).
\]
This rule first appears in print in a paper of Leibniz of October
1684~\cite{Leibniz1684}. The paper is written in Latin. It is the paper that
introduced the differential calculus.

Why did Leibniz need such a rule? The title of the paper names the problem. Leibniz wants a method for maxima,
minima and tangents \translated{which is not stopped by fractions or by
irrational quantities}{quae nec fractas nec irrationales quantitates
moratur}{Latin}. The
methods of his time first removed all square roots from an equation. Leibniz
describes this work as follows:
\begin{originalquote}{Latin}{calculo praestare sublatis irrationalibus, taediosissimae et aliquando insuperabilis operae}
to carry out by calculation after removing the irrationals, a most tedious
and sometimes insuperable labour
\end{originalquote}
\noindent\cite[p.~472]{Leibniz1684}. His answer is a set of rules. Each rule
acts on one symbol of a formula at a time. So a square root or a fraction
needs no preparation. He tests the rules on the law of refraction, which he
finds as a minimum. Then he remarks~\cite[p.~472]{Leibniz1684}:
\begin{originalquote}{Latin}{Cum alii doctissimi Viri multis ambagibus venati sint, quae hujus calculi peritus tribus lineis imposterum praestabit.}
Whereas other most learned men have hunted with many detours for what one
skilled in this calculus will henceforth achieve in three lines.
\end{originalquote}
\noindent The product rule is one of these rules. It gives the differential of
a product from the differentials of its factors. Nothing has to be expanded.

To define his differences, Leibniz draws a curve above a horizontal axis. A point $X$ on the axis has
abscissa $x$, and the curve has height $v$ above it. The tangent to the curve
at that point meets the axis at a point $B$. The segment $BX$ is called the
\emph{subtangent} \cite[p.~467]{Leibniz1684}. Leibniz chooses any segment $dx$. Then he defines $dv$ by
the proportion
\[
  dv : dx = v : BX .
\]
He calls $dv$ the \emph{differentia}, that is, the difference of $v$ \cite[p.~467]{Leibniz1684}.
Figure~\ref{fig:origins:subtangent} shows the construction.

\hbfigure{01-origins/leibniz-subtangent}{The construction of Leibniz \cite[p.~467]{Leibniz1684}. The tangent at $V$ meets the axis at $B$, and $BX$ is the subtangent. The difference $dv$ satisfies $dv : dx = v : BX$.}{fig:origins:subtangent}

With this construction, Leibniz states his rules. The first two rules say that constants have zero difference and that $d$ is
linear: \translated{Let $a$ be a given constant quantity; then $da$ will be
equal to $0$, and $d\,ax$ will be equal to $a\,dx$}{Sit $a$ quantitas data
constans, erit $da$ aequalis $0$, \& $d\,ax$ erit aequ.\ $a\,dx$}{Latin}. The
difference of a sum is the sum of the differences. The third rule is the
product rule: \translated{Multiplication: $d\,xv$ equals $x\,dv +
v\,dx$}{Multiplicatio, $d\,xv$ aequ.\ $x\,dv + v\,dx$}{Latin}%
~\cite[p.~467]{Leibniz1684}. The quotient rule follows on the same page.
Figure~\ref{fig:origins:product-area} gives the usual picture of the product
rule.

\hbfigure{01-origins/product-rule-area}{The product rule read as area. When $x$ and $v$ grow by $dx$ and $dv$, the rectangle $xv$ grows by $x\,dv + v\,dx$. The small corner $dx\,dv$ is of higher order.}{fig:origins:product-area}

Three features of this passage return in every later definition. First, the
rules are stated for symbols, and no limit appears. Leibniz says that $dx$ is
treated in the calculus like $y$ and $dy$, or like any other letter with its
differential~\cite[p.~467]{Leibniz1684}.

Second, constants are exactly the quantities that $d$ sends to zero. Dedekind
and Weber stress the same property two centuries later
(Section~\ref{sec:origins:function-fields}). Third, the product rule is the
only rule that links $d$ to multiplication. So a linear map with this rule
behaves like a derivative in every computation made of sums and products.

Today the third feature is the definition itself. Let $k$ be a commutative ring and $A$ a $k$-algebra. A \emph{derivation} of
$A$ is a $k$-linear map $D\colon A \to A$ with $D(xy) = D(x)y + xD(y)$ for all
$x, y \in A$ \cite[p.~206]{Jacobson1937}. If $A$ has a unit, then $D(1) = D(1\cdot 1) = 2D(1)$, so
$D(1) = 0$. By linearity, $D$ vanishes on the image of $k$. Induction on the
product rule gives the general Leibniz formula:
\[
  D^n(xy) = \sum_{i=0}^{n} \binom{n}{i} D^i(x)\, D^{n-i}(y).
\]
Jacobson uses this formula in 1937. It shows that in characteristic $p$ the
$p$-th power of a derivation is again a derivation
(Section~\ref{sec:origins:lie-theory}). Chapter~\ref{ch:basic-theory}
develops the definition and its first consequences.

For two centuries after Leibniz, the product rule stayed a rule of analysis.
It entered algebra through the theory of algebraic functions, as the next
section shows.
